\documentclass[11pt]{letter}

\usepackage{geometry}
\geometry{margin=1in}
\usepackage{setspace}
\usepackage{hyperref}

\signature{Decory Edwards}
\address{Decory Edwards \\ Johns Hopkins University \\ Department of Economics}

\begin{document}

\begin{letter}{Hiring Committee \\ Da Vinci}

\opening{Dear Hiring Committee,}

I am writing to apply for a Graduate Quant position at Da Vinci. I am currently a PhD candidate in Economics at Johns Hopkins University. My training combines mathematical reasoning, statistical modeling, and large-scale computational analysis. I am motivated by environments that reward rigorous thinking, careful implementation, and continuous improvement based on real performance. Da Vinci’s emphasis on collaboration, autonomy, and end-to-end ownership of ideas strongly aligns with how I approach quantitative work.

My research agenda is at the intersection of wealth inequality and macroeconomics. In my job market paper, I build and estimate a structural heterogeneous agent model with returns that differ across households. The core contribution is to show how heterogeneity in returns and income risk jointly shape the wealth distribution and consumption behavior. Relative to closely related work, my framework performs better at matching the lower and middle portions of the wealth distribution. This difference matters quantitatively. Models that focus primarily on returns heterogeneity can fit the top of the wealth distribution closely but tend to generate too much wealth among the bottom half of households. In my framework, income uncertainty generates a precautionary saving motive that produces genuinely low wealth households. This leads to a realistic distribution of marginal propensities to consume and aggregate consumption responses that align with empirical estimates.

A key feature of my research is the disciplined analysis of complex, noisy systems. I work with large datasets, design models to extract signal from uncertainty, and test performance through simulation and robustness checks. I am comfortable translating mathematical theory into practical algorithms and iterating based on empirical results. This approach maps directly to developing, monitoring, and improving trading strategies and research tools in live market environments.

I am particularly interested in Da Vinci because of its integrated approach across trading, research, and development. The opportunity to work closely with traders, researchers, and engineers to push ideas from analysis to production is especially appealing. I value environments where responsibility is shared, performance is measured objectively, and strong results are rewarded.

My economics training has provided a strong foundation in probability, optimization, and econometric reasoning. I am comfortable programming in Python and related scientific computing environments and enjoy working collaboratively to solve demanding problems under time pressure. I bring a proactive mindset, strong attention to detail, and a competitive but cooperative approach to team-based work.

I am enthusiastic about the opportunity to join Da Vinci’s team, including the prospect of relocation and training in Amsterdam. I see this role as an opportunity to apply my analytical training in a high-performance trading environment while continuing to grow as a quantitative researcher.

I would welcome the opportunity to contribute my skills and perspective to Da Vinci. Thank you for your time and consideration.

\closing{Sincerely,}

\end{letter}
\end{document}
